\section{Tupolev Tu-2}

The Tu-2 is a propeller-driven medium bomber that was developed during WWII and which entered service with Soviet forces in small numbers in 1942. Despite early successes, production was suspended in favor of fighters and the Pe-2 light bomber. Production restarted in 1943 of the improved Tu-2S version, and this saw service from early 1944 until the end of WW2. It also served with Chinese communist forces in the Chinese Civil War and with the PLAAF and KPAF the Korean War. Its NATO reporting name is Bat.

The PLAAF developed the Tu-2P to intercept overflights by ROCAF aircraft. The RP-1 radar from the J-5A (MiG-17PF) was installed in the nose and two 23 mm NR-23 cannons replaced the 20 mm cannons in the wings. Defensive armament was deleted. The radar was modified to eliminate the lower part of the scan, which reduced ground clutter at the cost of only being able to detect and track targets at the same of higher altitude.

ADCs are provided for:
\begin{adclist}
    \adcitem{Tu-2S}
    \adcitem{Tu-2P}
\end{adclist}